\section{Low-\texorpdfstring{$T_c$}{Tc} Regime: Mediator-Stabilized Composites}
  \label{sec:low-tc-regime}

  \subsection{Mediator-Induced Reduction of the Projective Gap}
    \label{subsec:mediator-gap}

    In conventional superconductors,
    pair formation is well described phenomenologically
    by an effective attractive interaction
    mediated by lattice vibrations.
    Within the present framework,
    this mechanism is reinterpreted
    as a reduction of the projective gap $\Delta_\Pi$
    associated with the $w=2$ composite class.

    The phonon field does not introduce
    a fundamentally new pairing degree of freedom.
    Rather, it modifies the local raccordement cost
    of the fiber phase between neighboring electronic excitations.
    The effective interaction lowers the minimal eigenvalue
    of the relaxation operator in Eq.~\ref{eq:projective-gap},
    thereby stabilizing the composite configuration.

    We denote the renormalized gap as
    \begin{equation}
      \Delta_\Pi^{\mathrm{eff}}
      =
      \Delta_\Pi^{(0)}
      -
      \delta \Delta_\Pi^{(\mathrm{ph})} ,
    \end{equation}
    where $\Delta_\Pi^{(0)}$ is the intrinsic stability cost
    in the absence of mediator coupling,
    and $\delta \Delta_\Pi^{(\mathrm{ph})}$
    encodes the phonon-induced reduction.

    A finite composite population becomes thermodynamically stable
    when
    \begin{equation}
      k_B T \lesssim \Delta_\Pi^{\mathrm{eff}} .
      \label{eq:lowTc-stability}
    \end{equation}
    Equation~\ref{eq:lowTc-stability}
    defines the onset of local composite formation,
    which coincides with the superconducting transition
    in weakly correlated materials.

  \subsection{Emergence of the BCS Gap Structure}
    \label{subsec:bcs-gap}

    Although the structural origin of pairing
    differs from the standard interaction-based description,
    the resulting excitation spectrum is equivalent.

    In the coherent regime,
    breaking a composite requires
    at least energy $2\Delta_\Pi^{\mathrm{eff}}$,
    corresponding to the creation of two $w=1$ excitations.
    The quasiparticle dispersion therefore takes the form
    \begin{equation}
      E(\mathbf{k})
      =
      \sqrt{
        \xi_{\mathbf{k}}^2
        +
        \Delta^2
      },
    \end{equation}
    with
    \begin{equation}
      \Delta
      \equiv
      \Delta_\Pi^{\mathrm{eff}} .
    \end{equation}
    The formal mapping between the projective gap
    $\Delta_\Pi$ and the quasiparticle gap
    $\Delta(\mathbf{k})$ measured by ARPES
    is derived in Appendix~\ref{app:mapping}.

    The conventional $s$-wave symmetry
    corresponds to the isotropic $w=2$ composite class
    that minimizes the raccordement cost
    in a lattice without strong anisotropic frustration.
    No additional angular structure is required,
    and the gap remains nodeless across the Fermi surface.

    The exponential dependence of $\Delta$
    on the effective coupling strength in BCS theory
    is reinterpreted here
    as the sensitivity of the minimal eigenvalue
    $\lambda_{\min}(\mathcal{L}_\Pi)$
    to weak mediator-induced renormalization
    of the fiber stability cost.

  \subsection{London Electrodynamics and Meissner Screening}
    \label{subsec:london-meissner}

    Once a macroscopic fraction of composites
    occupies a common phase sector,
    the free-energy functional
    of Eq.~\ref{eq:london-structure}
    governs long-wavelength electrodynamics.

    Varying the free energy with respect to $A$
    yields the London equation
    \begin{equation}
      \nabla^2 \mathbf{B}
      =
      \frac{1}{\lambda_L^2}
      \mathbf{B},
      \label{eq:london-equation}
    \end{equation}
    with penetration depth
    \begin{equation}
      \lambda_L^{-2}
      =
      \frac{4 e^2 \rho_s}{\hbar^2 m^*}.
    \end{equation}

    The magnetic field is therefore expelled
    from the bulk over the scale $\lambda_L$,
    demonstrating the Meissner effect.
    In this interpretation,
    the gauge field acquires an effective mass
    through phase locking of the fiber,
    analogous to the Anderson--Higgs mechanism
    in condensed matter systems.

  \subsection{Flux Quantization}
    \label{subsec:flux-quantization}

    A vortex corresponds to a topological defect
    in the fiber phase.
    Around a closed contour encircling the vortex core,
    \begin{equation}
      \oint \nabla \theta \cdot d\ell
      =
      2\pi n,
      \qquad n \in \mathbb{Z}.
    \end{equation}

    Imposing $D\theta = 0$ outside the core,
    one obtains
    \begin{equation}
      \oint A \cdot d\ell
      =
      \frac{\hbar}{2e}
      \oint \nabla \theta \cdot d\ell
      =
      n \frac{h}{2e}.
    \end{equation}

    The magnetic flux is therefore quantized in units
    \begin{equation}
      \Phi_0
      =
      \frac{h}{2e}.
    \end{equation}

    Flux quantization follows directly
    from the topological winding of the fiber phase
    and does not require an independent postulate.

  \subsection{Summary of the Low-\texorpdfstring{$T_c$}{Tc} Regime}
    \label{subsec:lowTc-summary}

    In the conventional regime,
    a mediator lowers the projective gap
    and stabilizes an isotropic $w=2$ composite.
    Global phase locking of these composites
    generates the London and Ginzburg--Landau structures,
    the Meissner effect,
    and flux quantization.

    The standard BCS phenomenology
    is therefore recovered as a special case
    of mediator-stabilized projective phase locking.
    The essential structural ingredient
    is not the specific form of the interaction,
    but the existence of a stable $w=2$ composite class
    and the energetic cost of breaking global fiber coherence.
